\section{参数来源与实际可用性}
\begin{center}\small
\begin{tabular}{p{.20\linewidth}p{.45\linewidth}p{.25\linewidth}}\toprule
对象 & 本版具体参数 & 依据\\\midrule
地理 & 源塔经纬度、完整道路弧、DEM & OSM/Copernicus记录\\
目标/导线 & 9角色，0.5m面，5样点，弧垂5m & LI23/KPLC迁移模板\\
几何修正 & 237个全景面向外2.01m & 原始模型几何矛盾修正\\
车辆 & 轴距2.5m，4×1.8×1.8m，6m走廊 & 项目参考底盘，未实测\\
地面冲突 & 车间净距0.5m；每候选1车 & 本版仿真设定\\
机体 & 7kg、速度5/12m/s、加速度3 & M300参考包络及设定\\
相机 & 5184×3888，6.83--119.94mm & H20边界＋理想针孔\\
验收代理 & GSD5mm/px、拖影1px、曝光1ms & 项目阈值，非缺陷检出率\\
能量 & 548Wh/对，初始100\%，最低20\% & TB60名义量＋假设\\
充电 & 852.444/182.667W，效率0.9 & 公开时长＋参考曲线\\
温度/冷却 & 0--40℃常容量，额外冷却0s & 未标定芯温动力学\\
天气 & 无雨、日间、风/阵风不超过10m/s & 参考运营门限\\
时间场适配 & 连续量插值；降水保留累计语义 & 消除阶跃风的姿态不连续\\
重规划 & 300s复核、10.96Wh偏差触发 & 明示仿真配置\\\bottomrule
\end{tabular}\end{center}
\begin{center}\small
\begin{tabular}{p{.29\linewidth}p{.61\linewidth}}\toprule
仍需绑定/验证 & 当前影响\\\midrule
日前日期与区域预报 & 本次数学定义未选择新的执行日；需冻结日期、覆盖及真实发布时间后构建数值实例\\
真实塔型/垂直基准 & 地图真值与模板/高程假设分开；不能声称实地精细数字孪生\\
道路/停车许可 & 拓扑可达不等于车辆转弯可行或允许占路\\
M300与充电硬件 & 仍无同机型实测功率、芯温和部分充电控制标定\\
237塔求解及执行 & 完整数学定义不代替SCIP转录、真实可行解与独立天气执行结果\\\bottomrule
\end{tabular}\end{center}
\clearpage

